\subsection{INVARIANT POLYNOMIAL FUNCTIONS}

Let \(\mathfrak{a}\) be a finite dimensional Lie algebra. In accordance with the conventions of no. 1, we identify the algebra \(\mathbf{S}(\mathfrak{a}^{*})\) , the algebra \(\mathbf{S}(\mathfrak{a})^{*gr}\) , and the algebra of polynomial functions on \(\mathfrak{a}\) . For all \(a \in \mathfrak{a}\) , let \(\theta(a)\) be the derivation of \(\mathbf{S}(\mathfrak{a})\) such that \(\theta(a)x = [a, x]\) for all \(x \in \mathfrak{a}\) . We know (Chap. I, §3, no. 2) that \(\theta\) is a representation of \(\mathfrak{a}\) on \(\mathbf{S}(\mathfrak{a})\) . Let \(\theta^{*}(\mathfrak{a})\) be the restriction of \(-^{t}\theta(a)\) to \(\mathbf{S}(\mathfrak{a}^{*})\) . Then \(\theta^{*}\) is a representation of \(\mathfrak{a}\) . If \(f \in \mathbf{S}^{n}(\mathfrak{a}^{*})\) , then \(\theta^{*}(a)f \in \mathbf{S}^{n}(\mathfrak{a}^{*})\) and, for \(x_{1}, \ldots, x_{n} \in \mathfrak{a}\) ,

\[
(\theta^ {*} (a) f) (x _ {1}, \dots , x _ {n}) = - \sum_ {1 \leq i \leq n} f (x _ {1}, \dots , x _ {i - 1}, [ a, x _ {i} ], x _ {i + 1}, \dots , x _ {n}).\tag{1}
\]

We deduce easily from (1) that \(\theta^{*}(a)\) is a derivation of \(\mathbf{S}(\mathfrak{a}^{*})\) . An element of \(\mathbf{S}(\mathfrak{a})\) (resp. \(\mathbf{S}(\mathfrak{a}^{*})\) ) that is invariant under the representation \(\theta\) (resp. \(\theta^{*}\) ) of a is called an invariant element of \(\mathbf{S}(\mathfrak{a})\) (resp. \(\mathbf{S}(\mathfrak{a}^{*})\) ).

Lemma 2. Let \(\rho\) be a finite dimensional linear representation of \(\mathfrak{a}\) , and \(m\) an integer \(\geq 0\) . The function \(x \mapsto \operatorname{Tr}(\rho(x)^m)\) on \(\mathfrak{a}\) is an invariant polynomial function.

Put \(g(x_{1},\ldots ,x_{m}) = \mathrm{Tr}(\rho (x_{1})\ldots \rho (x_{m}))\) for \(x_{1},\ldots ,x_{m}\in \mathfrak{a}\) . If \(x\in \mathfrak{a}\) , we have

\[
\begin{array}{l} - (\theta^ {*} (x) g) (x _ {1}, \ldots , x _ {m}) \\ \qquad = \sum_ {1 \leq i \leq m} \mathrm{Tr} (\rho (x _ {1}) \ldots \rho (x _ {i - 1}) [ \rho (x), \rho (x _ {i}) ] \rho (x _ {i + 1}) \ldots \rho (x _ {m})) \\ \qquad = \mathrm{Tr} (\rho (x) \rho (x _ {1}) \ldots \rho (x _ {m})) - \mathrm{Tr} (\rho (x _ {1}) \ldots \rho (x _ {m}) \rho (x)) = 0, \end{array}
\]

so \(\theta^{*}(x)g = 0\) . Let \(h\) be the symmetric multilinear form defined by

\[
h (x _ {1}, \dots , x _ {m}) = \frac {1}{m !} \sum_ {\sigma \in \mathfrak {S} _ {m}} g (x _ {\sigma (1)}, \dots , x _ {\sigma (m)}).
\]

For all \(x \in \mathfrak{a}\) , we have \(\theta^{*}(x)h = 0\) and \(\operatorname{Tr}(\rho(x)^m) = h(x, \ldots, x)\) , hence the lemma.

Lemma 3. Let E be a finite dimensional g-module, and \(x \in E\) . Then x is an invariant element of the g-module E if and only if \((\exp a_{\mathrm{E}}).x = x\) for every nilpotent element a of g.

The condition is clearly necessary. Assume now that it is satisfied. Let a be a nilpotent element of g. There exists an integer n such that \(a_{E}^{n}=0\) . For all \(t\in k\) , we have

\[
0 = \exp (t a _ {\mathrm{E}}). x - x = t a _ {\mathrm{E}} x + \frac {1}{2 !} t ^ {2} a _ {\mathrm{E}} ^ {2} x + \dots + \frac {1}{(n - 1) !} t ^ {n - 1} a _ {\mathrm{E}} ^ {n - 1} x,
\]

so \(a_{\mathrm{E}}x = 0\) . But the Lie algebra \(\mathfrak{g}\) is generated by its nilpotent elements (§4, no. 1, Prop. 1). Hence \(x\) is an invariant element of the \(\mathfrak{g}\) -module E. Q.E.D.

For any \(\xi \in \mathbf{GL}(\mathfrak{g})\) , let \(\mathbf{S}(\xi)\) be the automorphism of \(\mathbf{S}(\mathfrak{g})\) that extends \(\xi\) , and \(\mathbf{S}^*(\xi)\) the restriction to \(\mathbf{S}(\mathfrak{g}^*)\) of the contragredient automorphism of \(\mathbf{S}(\xi)\) . Then \(\mathbf{S}\) and \(\mathbf{S}^*\) are representations of \(\mathbf{GL}(\mathfrak{g})\) . If \(a\) is a nilpotent element of \(\mathfrak{g}\) , \(\theta(a)\) is locally nilpotent on \(\mathbf{S}(\mathfrak{g})\) and \(\mathbf{S}(\exp a) = \exp \theta(a)\) , so

\[
\mathbf {S} ^ {*} (\exp \operatorname{ad} a) = \exp \theta^ {*} (a).\tag{2}
\]

PROPOSITION 3. Let \(f\) be a polynomial function on \(\mathfrak{g}\) . The following conditions are equivalent:

\begin{enumerate}
\def\labelenumi{(\roman{enumi})}
\item
  \(f\circ s = f\) for all \(s\in \mathrm{Aut}_e(\mathfrak{g})\)
\item
  \(f\circ s = f\) for all \(s\in \mathrm{Aut}_0(\mathfrak{g})\)
\item
  \(f\) is invariant.
\end{enumerate}

The equivalence of (i) and (iii) follows from formula (2) and Lemma 3. It follows from this that (iii) implies (ii) by extension of the base field. The implication (ii) \(\Longrightarrow\) (i) is clear.

Note carefully that, if f satisfies the conditions of Prop. 3, f is not in general invariant under \(\operatorname{Aut}(\mathfrak{g})\) (Exerc. 1 and 2).

THEOREM 1. Let \(\mathrm{I}(\mathfrak{g}^*)\) be the algebra of invariant polynomial functions on \(\mathfrak{g}\) . Let \(i: \mathbf{S}(\mathfrak{g}^*) \to \mathbf{S}(\mathfrak{h}^*)\) be the restriction homomorphism.

\begin{enumerate}
\def\labelenumi{(\roman{enumi})}
\item
  The map \(i|\mathrm{I}(\mathfrak{g}^*)\) is an isomorphism from the algebra \(\mathrm{I}(\mathfrak{g}^*)\) to the algebra \(\mathbf{S}(\mathfrak{h}^*)^{\mathrm{W}}\) .
\item
  For any integer \(n \geq 0\) , let \(\mathrm{I}^{n}(\mathfrak{g}^{*})\) be the set of homogeneous elements of \(\mathrm{I}(\mathfrak{g}^{*})\) of degree n.~Then \(\mathrm{I}^{n}(\mathfrak{g}^{*})\) is the set of linear combinations of functions on g of the form \(x \mapsto \mathrm{Tr}(\rho(x)^{n})\) , where \(\rho\) is a finite dimensional linear representation of g.
\item
  Let \(l = \operatorname{rk}(\mathfrak{g})\) . There exist \(l\) algebraically independent homogeneous elements of \(\mathrm{I}(\mathfrak{g}^*)\) that generate the algebra \(\mathrm{I}(\mathfrak{g}^*)\) .
\end{enumerate}

\begin{enumerate}
\def\labelenumi{\alph{enumi})}
\item
  Let \(f \in \mathrm{I}(\mathfrak{g}^{*})\) and \(w \in W\) . There exists \(s \in \mathrm{Aut}_{e}(\mathfrak{g}, \mathfrak{h})\) such that \(s|_{\mathfrak{h}} = w\) (§2, no. 2, Cor. of Th. 2). Since f is invariant under s (Prop. 3), \(i(f)\) is invariant under w. Hence \(i(\mathrm{I}(\mathfrak{g}^{*})) \subset \mathbf{S}(\mathfrak{h}^{*})^{\mathrm{W}}\) .
\item
  We prove that, if \(f \in \mathrm{I}(\mathfrak{g}^{*})\) is such that \(i(f) = 0\) , then f = 0. Extending the base field if necessary, we can assume that k is algebraically closed. By Prop. 3, f vanishes on \(s(\mathfrak{h})\) for all \(s \in \operatorname{Aut}_{e}(\mathfrak{g})\) . Hence f vanishes on every Cartan subalgebra of g (Chap. VII, §3, no. 2, Th. 1), and in particular on the set of regular elements of g. But this set is dense in g for the Zariski topology (Chap. VII, §2, no. 2).
\item
  Let n be an integer \(\geq 0\) . Let \(L^{n}\) be the set of linear combinations of functions of the form \(x \mapsto \operatorname{Tr}(\rho(x)^{n})\) on g, where \(\rho\) is a finite dimensional linear representation of g. By Lemma 2, \(L^{n} \subset \operatorname{I}^{n}(\mathfrak{g}^{*})\) . Thus
\end{enumerate}

\[
i \left(\mathrm{L} ^ {n}\right) \subset i \left(\mathrm{I} ^ {n} \left(\mathfrak {g} ^ {*}\right)\right) \subset \mathbf {S} ^ {n} \left(\mathfrak {h} ^ {*}\right) ^ {\mathrm{W}}.
\]

By Cor. 2 of Prop. 2, \(\mathbf{S}^n (\mathfrak{h}^*)^{\mathrm{W}}\subset i(\mathrm{L}^n)\) . Hence \(i(\mathrm{I}^n (\mathfrak{g}^*)) = \mathbf{S}^n (\mathfrak{h}^*)^{\mathrm{W}}\) , which proves (i), and \(i(\mathrm{L}^n) = i(\mathrm{I}^n (\mathfrak{g}^*))\) so \(\mathrm{L}^n = \mathrm{I}^n (\mathfrak{g}^*)\) by \(b\) ). Thus (ii) is proved.

\begin{enumerate}
\def\labelenumi{\alph{enumi})}
\setcounter{enumi}{3}
\tightlist
\item
  Assertion (iii) follows from (i) and Chap. V, §5, no. 3, Th. 3.
\end{enumerate}

COROLLARY 1. Assume that \(\mathfrak{g}\) is simple. Let \(m_1, \ldots, m_l\) be the exponents of the Weyl group of \(\mathfrak{g}\) . There exist elements \(\mathrm{P}_1, \ldots, \mathrm{P}_l\) of \(\mathrm{I}(\mathfrak{g}^*)\) , homogeneous of degrees

\[
m _ {1} + 1, \dots , m _ {l} + 1,
\]

which are algebraically independent and generate the algebra \(\mathrm{I}(\mathfrak{g}^*)\) .

This follows from Th. 2 (i) and Chap. V, §6, no. 2, Prop. 3.

COROLLARY 2. Let B be a basis of R, \(R_{+}\) (resp. \(R_{-}\) ) the set of positive (resp. negative) roots of \((\mathfrak{g},\mathfrak{h})\) relative to B, \(n_{+}=\sum_{\alpha\in R_{+}}g^{\alpha}\) , \(n_{-}=\sum_{\alpha\in R_{-}}g^{\alpha}\) , \(\mathbf{S}(\mathfrak{h})\) the symmetric algebra of h, and J the ideal of \(\mathbf{S}(\mathfrak{g})\) generated by \(n_{+}\cup n_{-}\) .

\begin{enumerate}
\def\labelenumi{(\roman{enumi})}
\item
  \(\mathbf{S}(\mathfrak{g}) = \mathbf{S}(\mathfrak{h}) \oplus \mathbf{J}.\)
\item
  Let j be the homomorphism from the algebra \(\mathbf{S}(\mathfrak{g})\) to the algebra \(\mathbf{S}(\mathfrak{h})\) defined by the preceding decomposition of \(\mathbf{S}(\mathfrak{g})\) . Let \(\mathrm{I}(\mathfrak{g})\) be the set of invariant elements of \(\mathbf{S}(\mathfrak{g})\) . Let \(\mathbf{S}(\mathfrak{h})^{\mathrm{W}}\) be the set of elements of \(\mathbf{S}(\mathfrak{h})\) invariant under the operation of W. Then \(j|I(\mathfrak{g})\) is an isomorphism from \(\mathrm{I}(\mathfrak{g})\) to \(\mathbf{S}(\mathfrak{h})^{\mathrm{W}}\) .
\end{enumerate}

Assertion (i) is clear. The Killing form defines an isomorphism from the vector space \(\mathfrak{g}^*\) to the vector space \(\mathfrak{g}\) , which extends to an isomorphism \(\xi\) from the \(\mathfrak{g}\) -module \(\mathbf{S}(\mathfrak{g}^*)\) to the \(\mathfrak{g}\) -module \(\mathbf{S}(\mathfrak{g})\) . We have \(\xi (\mathrm{I}(\mathfrak{g}^*)) = \mathrm{I}(\mathfrak{g})\) . The orthogonal complement of \(\mathfrak{h}\) with respect to the Killing form is \(\mathfrak{n}_{+} + \mathfrak{n}_{-}\) (§2, no. 2, Prop. 1). If we identify \(\mathfrak{h}^*\) with the orthogonal complement of \(\mathfrak{n}_{+} + \mathfrak{n}_{-}\) in \(\mathfrak{g}^*\) , then \(\xi (\mathfrak{h}^*) = \mathfrak{h}\) , so \(\xi (\mathbf{S}(\mathfrak{h}^*)) = \mathbf{S}(\mathfrak{h})\) and \(\xi (\mathbf{S}(\mathfrak{h}^*)^{\mathrm{W}}) = \mathbf{S}(\mathfrak{h})^{\mathrm{W}}\) . Finally, \(\xi^{-1}(\mathrm{J})\) is the set of polynomial functions on \(\mathfrak{g}\) that vanish on \(\mathfrak{h}\) . This proves that \(\xi\) transforms the homomorphism \(i\) of Th. 1 into the homomorphism \(j\) of Cor. 2. Thus assertion (ii) follows from Th. 1 (i).

PROPOSITION 4. Let \(\mathfrak{a}\) be a semi-simple Lie algebra, \(l\) its rank. Let I (resp. I') be the set of elements of \(\mathbf{S}(\mathfrak{a}^{*})\) (resp. \(\mathbf{S}(\mathfrak{a})\) ) invariant under the representation induced by the adjoint representation of \(\mathfrak{a}\) . Let Z be the centre of the enveloping algebra of \(\mathfrak{a}\) .

\begin{enumerate}
\def\labelenumi{(\roman{enumi})}
\item
  I and \(\mathrm{I}'\) are graded polynomial algebras (Chap. V, §5, no. 1) of transcendance degree \(l\) .
\item
  \(Z\) is isomorphic to the algebra of polynomials in \(l\) indeterminates over \(k\) .
\end{enumerate}

The canonical filtration of the enveloping algebra of \(\mathfrak{a}\) induces a filtration of \(Z\) . By Chap. I, §2, no. 7, Th. 1 and p.~25, gr Z is isomorphic to I'. In view of Commutative Algebra, Chap. III, §2, no. 9, Prop. 10, it follows that (i) \(\Longrightarrow\) (ii).

On the other hand, Th. 1 and its Cor. 2 show that (i) is true whenever \(\mathfrak{a}\) is split. The general case reduces to that case in view of the following lemma:

Lemma 4. \(^{2}\) Let \(A = \bigoplus_{n \geq 0} A^{n}\) be a graded k-algebra, \(k'\) an extension of k, and \(A' = A \otimes_{k} k'\) . Assume that \(A'\) is a graded polynomial algebra over \(k'\) . Then A is a graded polynomial algebra over k.

We have \(A^{\prime0}=k'\) , so \(A^{0}=k\) . Put \(A_{+}=\bigoplus_{n\geq1}A^{n}\) and \(P=A_{+}/A_{+}^{2}\) . Then P is a graded vector space, and there is a graded linear map \(f:P\to A_{+}\) of degree zero such that the composite with the canonical projection \(A_{+} \rightarrow P\) is the identity on P. Give \(\mathbf{S}(P)\) the graded structure induced by that of P (Algebra, Chap. III, p.~506). The homomorphism of k-algebras \(g : \mathbf{S}(P) \rightarrow A\) that extends f (Algebra, Chap. III, p.~497) is a graded homomorphism of degree 0; an immediate induction on the degree shows that g is surjective.

Lemma 5. A is a graded polynomial algebra if and only if P is finite dimensional and g is bijective.

If P is finite dimensional, \(\mathbf{S}(\mathrm{P})\) is clearly a graded polynomial algebra, and so is A if g is bijective. Conversely, assume that A is generated by algebraically independent homogeneous elements \(x_{1},\ldots,x_{m}\) of degrees \(d_{1},\ldots,d_{m}\) . Let \(\bar{x}_{i}\) be the image of \(x_{i}\) in P. It is immediate that the \(\bar{x}_{i}\) form a basis of P; since \(\bar{x}_{i}\) is of degree \(d_{i}\) , it follows that \(\mathbf{S}(\mathrm{P})\) and A are isomorphic; in particular, \(\dim\mathbf{S}(\mathrm{P})^{n}=\dim\mathrm{A}^{n}\) for all n.~Since g is surjective, it is necessarily bijective.

Lemma 4 is now immediate. Indeed, Lemma 5, applied to the \(k'\) -algebra \(\mathrm{A}'\) , shows that \(g \otimes 1: \mathbf{S}(\mathrm{P}) \otimes k' \to \mathrm{A} \otimes k'\) is bijective, and hence so is \(g\) .

PROPOSITION 5. We retain the notations of Prop. 4, and denote by p the ideal of \(\mathbf{S}(\mathfrak{a}^{*})\) generated by the homogeneous elements of I of degree \(\geq 1\) . Let \(x \in a\) . Then x is nilpotent if and only if \(f(x) = 0\) for all \(x \in p\) .³

Extending the base field if necessary, we can assume that \(\mathfrak{a} = \mathfrak{g}\) is splittable. Assume that \(x\) is nilpotent. For any finite dimensional linear representation \(\rho\) of \(\mathfrak{g}\) , and any integer \(n \geq 1\) , we have \(\operatorname{Tr}(\rho(x)^n) = 0\) , so \(f(x) = 0\) for all homogeneous \(f \in \mathrm{I}(\mathfrak{g}^*)\) of degree \(\geq 1\) (Th. 1 (ii)), and hence \(f(x) = 0\) for all \(f \in \mathfrak{p}\) . Conversely, if \(f(x) = 0\) for all \(f \in \mathfrak{p}\) , then \(\operatorname{Tr}((\operatorname{ad} x)^n) = 0\) for all \(n \geq 1\) (Th. 1 (ii)), so \(x\) is nilpotent.

*Remarks 1) Let \(\mathrm{P}_1, \ldots, \mathrm{P}_l\) be algebraically independent homogeneous elements of I that generate the algebra I. Then \((\mathrm{P}_1, \ldots, \mathrm{P}_l)\) is an \(\mathbf{S}(\mathfrak{a}^*)\) -regular sequence (Chap. V, §5, no. 5). Indeed, extending the base field if necessary, we can assume that \(\mathfrak{a} = \mathfrak{g}\) is splittable. Now let \(\mathrm{N} = \dim \mathfrak{g}\) , and let

\[
\left(\mathrm{Q} _ {1}, \dots , \mathrm{Q} _ {\mathrm{N} - l}\right)
\]

be a basis of the orthogonal complement of \(\mathfrak{h}\) in \(\mathfrak{g}^*\) . Let \(\mathfrak{m}\) be the ideal of \(\mathbf{S}(\mathfrak{g}^*)\) generated by \(\mathrm{P}_1, \ldots, \mathrm{P}_l, \mathrm{Q}_1, \ldots, \mathrm{Q}_{\mathrm{N}-l}\) . Then \(\mathbf{S}(\mathfrak{g}^*)\mathfrak{m}\) is isomorphic to \(\mathbf{S}(\mathfrak{h}^*)/\mathrm{J}\) , where \(\mathrm{J}\) is the ideal of \(\mathbf{S}(\mathfrak{h}^*)\) generated by \(i(\mathrm{P}_1), \ldots, i(\mathrm{P}_l)\) . By Th. 1 and Chap. V, §5, no. 2, Th. 2, \(\mathbf{S}(\mathfrak{h}^*)/\mathrm{J}\) is a finite dimensional vector space, and hence so is \(\mathbf{S}(\mathfrak{g}^*)/\mathfrak{m}\) . By a result of Commutative Algebra, it follows that \((\mathrm{P}_1, \ldots, \mathrm{P}_l, \mathrm{Q}_1, \ldots, \mathrm{Q}_{\mathrm{N}-l})\) is an \(\mathbf{S}(\mathfrak{g}^*)\) -regular sequence, and a fortiori so is \((\mathrm{P}_1, \ldots, \mathrm{P}_l)\) .

\begin{enumerate}
\def\labelenumi{\arabic{enumi})}
\setcounter{enumi}{1}
\tightlist
\item
  The algebra \(\mathbf{S}(\mathfrak{a}^*)\) is a graded free module over I. Indeed, this follows from Prop. 4, Remark 1, and Chap. V, §5, no. 5, Lemma 5.*
\end{enumerate}

